\cxtitle{Ryan O'Donnell's matrix conjecture}

\begin{cxcredits}
\posedby{R.~O'Donnell and J.~Wright \citeyearpar{ODonnell2015AlmostDiagonal}}
\foundby{GPT-5.6 Pro}{2026-08}
\formalizedby{GPT-5.6 Pro}
\auditedby{GPT-5.6 Pro}
\contributedby{S.~Sra}
\end{cxcredits}

\begin{cxcontext}
For a Hermitian matrix $X$ write $\lVert X\rVert_1=\operatorname{tr}|X|$ for its trace norm.  For $R\in\mathbb C^{n\times n}$, let $D=\operatorname{Diag}(R)=\operatorname{diag}(d_1,\ldots,d_n)$ be its diagonal
part, and let $\Lambda(R)$ be the diagonal matrix carrying the eigenvalues
$\lambda_1\geq\cdots\geq\lambda_n$ in sorted order, so that
\[
  \epsilon:=\lVert D-\Lambda(R)\rVert_1=\nlsum_i|\lambda_i-d_{ii}|.
\]
The source states its conclusion asymptotically, as $\lVert R-D\rVert_1=O(\sqrt\epsilon)$; an
absolute constant in that $O(\cdot)$ is exactly an absolute $c$ with
$\lVert R-D\rVert_1^2\leq c\,\epsilon$, which is the form refuted below.

% Neither $\epsilon$ nor the source's conclusion specifies how the diagonal
% entries are to be ordered against the sorted spectrum.  In the family below
% both $(\lambda_i)$ and $(d_i)$ are strictly decreasing in the displayed index
% order, and the common sorted order minimises $\sum_i|\lambda_i-d_{\sigma(i)}|$
% over permutations $\sigma$; the witness therefore makes $\epsilon$ as small
% as any matching allows, and the refutation does not rely on a favorable
% pairing.
\end{cxcontext}

\begin{cxsource}{odonnell-matrix-conjecture}
R.~O'Donnell, answering a question of Nengkun Yu on MathOverflow
\cite{ODonnell2015AlmostDiagonal}, reports that he and J.~Wright had been
considering the problem in connection with quantum tomography, and closes by
conjecturing the unit-trace case.  The same conjecture was put to S.~Sra by
O'Donnell in personal communication.
\end{cxsource}

\begin{cxstatement}{odonnell-matrix-conjecture}
The question answered was whether a positive semidefinite $M$ with
$\lVert D(M)-\Lambda(M)\rVert_1\leq\epsilon$ must lie within $2\epsilon$ of
% some diagonal matrix.  O'Donnell observes that for the variant he and Wright
% cared about --- $M$ close to \emph{its own} diagonal --- the answer is no in
% general, and then
conjectures the normalized case:

\begin{quote}
``If you restrict $M$ to have trace $1$, then we felt the desired result was
true but with a weaker conclusion of $O(\sqrt\epsilon)$ rather than
$2\epsilon$.''
\end{quote}

Written out, and with the implied constant absolute: for every unit-trace
positive semidefinite $R$ with eigenvalues $\lambda_i$ and diagonal entries
$d_{ii}$,
\[
  \lVert R-\operatorname{Diag}(R)\rVert_1^2
  \leq c\nlsum_i |\lambda_i-d_{ii}|
\]
for some absolute constant $c$ independent of $n$.
\end{cxstatement}

\begin{cxsummary}{odonnell-matrix-conjecture}
Every absolute $c$ fails: an exact real-symmetric density-matrix family has ratio greater than $(\log_2 n)/32$.
\end{cxsummary}

\begin{cxcertificate}{odonnell-matrix-conjecture}
An exact algebraic Givens-rotation family with a rational dyadic lower bound, accompanied by a standard-library exact-arithmetic verifier.
\end{cxcertificate}

\begin{cxrefutation}
The conjecture \cite{ODonnell2015AlmostDiagonal} asks whether the off-diagonal
trace norm of a density matrix can be controlled, up to an absolute constant,
by the square root of the $\ell_1$ displacement between its spectrum and its
diagonal.  The following exact family shows logarithmic divergence.

\begin{theorem}[\statusfalse: O'Donnell's matrix conjecture]\label{thm:odonnell-matrix-conjecture}
For every $C>0$ there exist an integer $n$ and a real symmetric positive
semidefinite matrix $R\in\mathbb R^{n\times n}$ with
$\operatorname{tr}R=1$ such that, writing
$D=\operatorname{Diag}(R)=\operatorname{diag}(d_1,\ldots,d_n)$ and listing the
eigenvalues $\lambda_1\geq\cdots\geq\lambda_n$,
\[
  \frac{\lVert R-D\rVert_1^2}
       {\sum_{i=1}^n |\lambda_i-d_i|}>C.
\]
Moreover, both $(\lambda_i)$ and $(d_i)$ are strictly decreasing.
\end{theorem}

\begin{proof}
Fix an integer $m\geq2$ and put
\[
  n=2^m,\qquad q=n,\qquad
  H=H_{n-1}=\sum_{i=1}^{n-1}\frac1i,
\]
\[
  B=1+q^2H,\qquad
  \delta=\frac1{2B},\qquad
  \eta=\frac1{2n}.
\]
Define positive spectral gaps by
\[
  g_1=\delta(1+q^2),
  \qquad
  g_i=\delta\frac{q^2}{i^2}
  \quad (2\leq i\leq n-1),
\]
and set
\[
  \lambda_n=\eta,
  \qquad
  \lambda_i=\eta+\sum_{k=i}^{n-1}g_k
  \quad (1\leq i\leq n-1).
\]
Then $\lambda_1>\cdots>\lambda_n>0$, and
\[
\begin{aligned}
  \sum_{i=1}^n\lambda_i
  &=n\eta+\sum_{i=1}^{n-1}i g_i \\
  &=\frac12+\delta\left(1+q^2+q^2\sum_{i=2}^{n-1}\frac1i\right)
    =\frac12+\delta B=1.
\end{aligned}
\]
Thus $\Lambda=\operatorname{diag}(\lambda_1,\ldots,\lambda_n)$ is a density
matrix.

For $1\leq i\leq n-1$, put
\[
  r_i=\frac qi,
  \qquad
  c_i=\frac{q}{\sqrt{q^2+i^2}},
  \qquad
  s_i=\frac{i}{\sqrt{q^2+i^2}}.
\]
Hence $c_i^2+s_i^2=1$ and $c_i/s_i=r_i$.  Let $G_i$ be the identity except
for the following orthogonal block in coordinates $i,i+1$:
\[
  \begin{pmatrix}
    c_i&-s_i\\
    s_i& c_i
  \end{pmatrix}.
\]
Starting from $R^{(0)}=\Lambda$, define successively
\[
  R^{(i)}=G_iR^{(i-1)}G_i^{\mathsf T},
  \qquad 1\leq i\leq n-1,
\]
and let $R=R^{(n-1)}$.  This is an exact algebraic matrix.  Since it is
orthogonally similar to $\Lambda$, it is real symmetric, positive semidefinite,
has trace one, and has eigenvalues $(\lambda_i)$.

We next compute its diagonal exactly.  Before the first rotation, the active
diagonal gap is
\[
  \lambda_1-\lambda_2=g_1=\delta(1+r_1^2).
\]
Before rotation $G_i$ for $i\geq2$, coordinate $i$ has diagonal entry
$\lambda_i+\delta$, coordinate $i+1$ has diagonal entry $\lambda_{i+1}$,
and coordinate $i+1$ has not previously been touched.  Thus the active
$2\times2$ principal block is diagonal and its diagonal gap is
\[
  (\lambda_i+\delta)-\lambda_{i+1}
  =g_i+\delta
  =\delta(1+r_i^2).
\]
In both cases,
\[
  s_i^2\,\delta(1+r_i^2)=\delta.
\]
Consequently, $G_i$ transfers exactly $\delta$ from diagonal position $i$ to
position $i+1$.  Induction gives
\[
  d_1=\lambda_1-\delta,
  \qquad
  d_i=\lambda_i\quad(2\leq i\leq n-1),
  \qquad
  d_n=\lambda_n+\delta.
\]
These entries are strictly decreasing: at the first edge the remaining gap is
$g_1-\delta=\delta q^2>0$, at every interior edge it is $g_i>0$, and at the
last edge it is
\[
  g_{n-1}-\delta
  =\delta\left(\frac{q^2}{(n-1)^2}-1\right)>0
\]
because $q=n$.  Therefore
\[
  \sum_{i=1}^n|\lambda_i-d_i|=2\delta.
\]

The same rotation calculation gives useful exact off-diagonal entries.  At
step $i$, the newly created $(i,i+1)$ entry has magnitude
\[
  c_is_i\,\delta(1+r_i^2)=\delta r_i.
\]
For $i\leq n-2$, the next rotation multiplies this entry by $c_{i+1}$; all
later rotations leave it unchanged.  Hence
\[
  |R_{i,i+1}|=
  \begin{cases}
    \delta r_i c_{i+1},&1\leq i\leq n-2,\\
    \delta r_{n-1},&i=n-1.
  \end{cases}
\]
Let $A=R-D$.  Pinch $A$ to the disjoint coordinate blocks
$(1,2),(3,4),\ldots,(n-1,n)$.  Trace norm is contractive under pinching, and
each retained zero-diagonal $2\times2$ block has trace norm twice the magnitude
of its off-diagonal entry.  If
\[
  O=\sum_{\substack{1\leq i\leq n-1\\ i\ \mathrm{odd}}}\frac1i,
\]
then $q=n$ implies
\[
  c_j^2=\frac{q^2}{q^2+j^2}>\frac12
  \qquad(1\leq j\leq n-1).
\]
Using the multiplier $1>1/\sqrt2$ for the final edge as well, we obtain
\[
  \lVert A\rVert_1
  \geq 2\sum_{\substack{1\leq i\leq n-1\\ i\ \mathrm{odd}}}|R_{i,i+1}|
  >\sqrt2\,\delta q O.
\]
It follows that
\[
\begin{aligned}
  \frac{\lVert R-D\rVert_1^2}
       {\sum_i|\lambda_i-d_i|}
  &>\frac{2\delta^2q^2O^2}{2\delta}
    =\delta q^2O^2
    =\frac{q^2O^2}{2(1+q^2H)}\\
  &>\frac{O^2}{2(H+1)}.
\end{aligned}
\]

It remains only to estimate the two harmonic sums.  In each dyadic block
$[2^k,2^{k+1})$, for $1\leq k\leq m-1$, there are $2^{k-1}$ odd integers,
each contributing more than $2^{-(k+1)}$.  Therefore
\[
  O>1+\frac{m-1}{4}=\frac{m+3}{4}.
\]
Also
\[
  H_{n-1}<1+\log(n-1)<m+1.
\]
Consequently,
\[
  \frac{\lVert R-D\rVert_1^2}
       {\sum_i|\lambda_i-d_i|}
  >\frac{(m+3)^2}{32(m+2)}
  >\frac{m}{32}.
\]
Given any real $C>0$, choose an integer $m>32C$.  The resulting exact matrix
has ratio greater than $C$, proving that no dimension-independent constant can
exist.
\end{proof}

\begin{remark}[Scope]
What falls is the unit-trace conjecture of \cite{ODonnell2015AlmostDiagonal},
namely $\lVert R-\operatorname{Diag}(R)\rVert_1=O(\sqrt\epsilon)$ with an
absolute implied constant.  Two neighbouring statements are untouched.  The
question originally posed there, whether such an $M$ lies within $2\epsilon$ of
\emph{some} diagonal matrix, is not the target; nor is the unnormalised variant,
which the same answer already resolves in the negative by a perturbative
argument.  The witness happens to settle the own-diagonal form of both, since
$\lVert R-D\rVert_1/\epsilon$ diverges along the family, but only the unit-trace
conjecture was posed as plausible and only it is claimed here.
\end{remark}
\end{cxrefutation}


%%% Local Variables:
%%% mode: LaTeX
%%% TeX-master: "../../tex/main"
%%% End:
